\section{Research Plan}
\label{sec:research-plan}

The goal of this proposal, restated from Section~\ref{sec:research-introduction},
is to comprehensively characterize how well current VLMs perform the core CV
tasks of disaster response -- detection, segmentation, and classification of
flooded areas, burned areas, and damaged roads and infrastructure (including
buildings and debris) -- once evaluation is extended across the fuller
landscape of disaster RS datasets and stratified by imagery provider and
sensor, and to establish whether a consolidated benchmark or fine-tuned model
can close the gaps found. A key targeted contribution of this work would be the modular framework of evaluation allowing us
and the wider community to plugin new datasets or models and running evaluation out of the box with central dashboards of performance.

\subsection{Research questions}

\textbf{RQ1 (Is the ``cross-sensor gap'' actually a cross-provider or
cross-resolution gap?)} Re-evaluate DisasterM3's 14 benchmarked and 4
fine-tuned VLMs with results disaggregated by imagery \emph{provider} (Maxar
vs.\ Capella vs.\ Umbra) and by \emph{resolution regime}, not only by the
optical/SAR modality split the paper reports. If accuracy tracks provider or
resolution more tightly than it tracks modality, the ``cross-sensor gap'' is
better described (and better fixed) as a provider- or resolution-robustness
problem.

\textbf{RQ2 (Does performance generalize across datasets, or is it
DisasterM3-specific?)} Extend the same stratified evaluation to MONITRS
(event-level temporal/textual narrative QA) and to a curated, license-cleared
subset of the datasets catalogued in our survey
(Section~\ref{sec:dataset-landscape}) -- e.g.\ BRIGHT and xBD for building
damage, Sen1Floods11/WorldFloods/C2S-MS Floods for flood segmentation on
Sentinel imagery, FloodNet and RescueNet for UAV-scale imagery, CanadaFire2023
and Sen2Fire for burned-area segmentation reformatted into a common
instruction/prompt schema. This tests whether DisasterM3-fine-tuned VLMs
generalize to providers, sensors, and hazards they were not tuned on, or
whether they have overfit to DisasterM3's specific Maxar/Capella/Umbra
distribution. Most likely, the already low performance values shown in DisasterM3, will be even lower which would make a compelling case
to develop an own model.

\textbf{RQ3 (Where are the task$\times$hazard coverage gaps?)} Map the
detection/segmentation/classification $\times$ \{flood, burned area, damaged
road, damaged building, debris\} grid the user community actually needs
against what DisasterM3's nine tasks cover today. DisasterM3 already supplies
damaged-building counting and damaged-road estimation; burned-area
segmentation and debris-specific detection are not represented in its task
suite and would need to be sourced from other survey entries (e.g.\
CanadaFire2023, Satellite Burned Area Dataset, RescueNet, ISBDA) and
reformatted into the same instruction schema.

\textbf{RQ4 (Can a consolidated model close the gaps found in RQ1--3?)} Once
RQ1--3 have identified where and how badly current VLMs fail, fine-tune one
or two open VLMs (e.g.\ Qwen2.5-VL, InternVL3, following DisasterM3's own
choice of base models) on a merged instruction corpus that is explicitly
balanced across providers, resolution regimes, and task$\times$hazard cells,
and test whether this closes the provider/resolution gaps that a DisasterM3-
only fine-tune does not.

\subsection{Approach}

\begin{enumerate}
  \item \textbf{Stratified re-evaluation harness (RQ1).} Build an evaluation
  harness around DisasterM3's own released test set and model checkpoints,
  adding provider and resolution metadata (already implicit in DisasterM3's
  data but not surfaced in its reported metrics) so every existing result can
  be re-sliced without any new training. This is the fastest step to execute
  and, on its own, already answers RQ1.
  \item \textbf{Dataset intake and harmonization (RQ2--3).} From the 44-dataset
  survey, filter to datasets that are downloadable and permissively licensed,
  prioritizing coverage of task$\times$hazard cells and provider/resolution
  regimes that DisasterM3 and BRIGHT under-represent (Sentinel-class medium
  resolution, UAV-scale imagery, burned area, debris). Convert their existing
  labels (masks, boxes, class labels) into instruction-style prompts
  comparable in format to DisasterM3's, and fold in MONITRS's event-level QA
  as a separate, temporally-oriented evaluation track rather than forcing it
  into the same per-image schema.
  \item \textbf{Extended stratified benchmarking (RQ2--3).} Run the same 14
  VLMs DisasterM3 benchmarks, plus its 4 fine-tuned checkpoints (if published), as well as newer models (for e.g. we don't see EarthMind in the eval), through the
  harmonized harness, reporting every metric stratified by provider,
  resolution regime, hazard type, and task -- producing, to our knowledge,
  the first provider/resolution-stratified view of VLM performance on
  disaster RS imagery.
  \item \textbf{Consolidated fine-tuning (RQ4, contingent on Steps 1--3).}
  If Steps 1--3 confirm systematic, provider- or resolution-linked failure
  modes (rather than idiosyncratic per-dataset noise), construct a
  provider/task-balanced instruction-tuning corpus from the harmonized data
  and fine-tune 1--2 open VLMs, evaluating them through the same stratified
  harness to test whether the gaps close.
\end{enumerate}

The stratified benchmarking harness and results (Steps 1--3) constitute a
complete, releasable contribution independent of whether Step 4 is pursued,
which keeps the proposal viable even if fine-tuning compute or licensing
constraints limit what Step 4 can cover.

\subsection{Risks and open questions}

We haven't found DisasterM3's fine-tuned models online. A further check is needed but perhaps directly contacting the authors could help.
Furthermore, the license and redistribution terms of
each survey dataset must be checked individually before inclusion in Step 2;
several entries in the survey record license as unspecified or unverified.
Reformatting mask/box-labeled datasets into instruction prompts (Step 2) is
itself a modeling choice that can introduce its own bias (e.g.\ how a
segmentation mask is turned into a natural-language question), and should be
validated.